\documentclass[11pt]{article}

\usepackage[utf8]{inputenc}
\usepackage[T1]{fontenc}
\usepackage{lmodern}
\usepackage{geometry}
\usepackage{amsmath,amssymb,amsfonts,amsthm,mathtools}
\usepackage{bm}
\usepackage{enumitem}
\usepackage{microtype}
\usepackage{hyperref}
\usepackage{titlesec}
\usepackage{booktabs}
\usepackage{array}
\usepackage{setspace}
\usepackage{mathrsfs}
\usepackage{physics}

\geometry{margin=1in}
\hypersetup{colorlinks=true, linkcolor=black, urlcolor=blue, citecolor=black}
\setlist[enumerate]{topsep=0.4em,itemsep=0.35em}
\setlist[itemize]{topsep=0.3em,itemsep=0.25em}
\titleformat{\section}{\large\bfseries}{\thesection.}{0.5em}{}
\titleformat{\subsection}{\normalsize\bfseries}{\thesubsection.}{0.5em}{}
\setstretch{1.06}

\newcommand{\Aether}{\AE{}ther}
\newcommand{\Aflow}{\AE{}ther-flow}
\newcommand{\TheoryName}{\Aflow{} Interpretation of Relativity}
\newcommand{\TheoryProgram}{\Aether{} / \Aflow{} framework}
\newcommand{\CompletedMetric}{g^{\sharp}}
\newcommand{\FoundationCore}{\mathfrak F^{\mathrm{fdn,core}}}

\newtheorem{definition}{Definition}
\newtheorem{proposition}{Proposition}
\newtheorem{remark}{Remark}
\newtheorem{corollary}{Corollary}
\newtheorem{theorem}{Theorem}

\title{\textbf{The \TheoryName{}}\\[0.4em]
\large Primitive-Reservoir Observer-Reduction\\
\large Foundation Frontier Next Benchmark Criterion}
\author{}
\date{}

\begin{document}

\maketitle

\begin{abstract}
The active primitive-reservoir observer-reduction line now sits at a more disciplined routing point than another same-output deeper-origin relay. The recorded ground-from-foundation continuation already moved the open burden to foundation scope, and the ground-generating substrate foundation dynamics necessity / uniqueness theorem then spent the honest benchmark-facing compression move available there: the benchmark-relevant foundation core is forced up to core-faithful equivalence once the fixed ground package is required. After that result, another shell, lift, support, reservoir, or ground restatement no longer counts as the right next benchmark-facing manuscript, and neither does a reopening of stationary reduction, stationary Hessians, spectral generators, or selector mechanics.

This note records the resulting criterion. The deeper foundation-from-root theorem remains a live manuscript and a legitimate deeper continuation, but under the current routing safeguard it does not by itself supersede the benchmark-facing status of the foundation-core compression theorem. The next benchmark-facing manuscript must therefore do one of two sharper things: either derive or justify the ground-generating substrate foundation dynamics from still more primitive substrate structure in a way that changes benchmark-facing status rather than merely relocating the unexplained layer deeper, or prove a genuinely smaller collapse strictly inside the benchmark-relevant foundation core \(\FoundationCore\).

The benchmark boundary remains fixed throughout. The one operative metric is still exactly \(\CompletedMetric\), the ordinary matter route is unchanged, there is no second operative metric, the low-energy matter law is not enlarged, and no stronger promotion language is licensed. The positive result of the note is therefore routing and scope discipline: it fixes what the next honest benchmark-facing gain at foundation scope is allowed to be.
\end{abstract}

\section{Introduction}

The active derivational program of \TheoryName{} is no longer at a stage where depth alone counts as benchmark-facing progress. The benchmark package is already fixed, the same-package observer-side remainder gate has already been settled on the active completed branch, and the upstream compression safeguard now blocks recursive depth drift. Once a manuscript preserves the same benchmark one-metric output while merely pushing the unexplained substrate layer one step deeper, it must supply a genuine new routing gain before it can honestly count as the default next move.

At current scope, the ground-generating substrate foundation dynamics necessity / uniqueness theorem is the manuscript that supplies that routing gain. Its theorem-level content is not another restatement of lower layers. It is a benchmark-relevant compression verdict at foundation scope: within the recorded foundation class, the core data that matter for the active line are no longer free once the fixed ground package is demanded. That is the live benchmark-facing frontier.

The deeper foundation-from-root theorem remains useful and active, but the present note records why it does not change that frontier on its own. The repository therefore needs a short benchmark criterion so that future work is routed honestly and the project does not drift into repeated same-output deeper-origin relays.

\paragraph{Framework and claim boundary.}
\TheoryProgram{} denotes the broader \Aether{} / \Aflow{} framework built on the claims that the \Aether{} is the underlying four-dimensional substrate of reality, the \Aflow{} is its intrinsic ordered motion, observed three-dimensional space is the local experiential slice of that deeper substrate, S-time is the experienced order of change arising from matter, light, and the \Aflow{}, and observed expansion is the three-dimensional appearance of deeper four-dimensional ordered motion. Within that framework, \TheoryName{} denotes the exact-closure theory in which the effective gravitational dynamics are adopted to be exactly Einsteinian with universal matter coupling. In the active sequence, \emph{adoption} means use of that established relativistic dynamics without claiming substrate derivation, while \emph{derivation} is reserved for a first-principles recovery from explicit substrate variables.

The present note stays strictly inside that boundary. It does not present a completed first-principles derivation of gravity. It records only which kind of next manuscript would now count as an honest benchmark-facing gain at the current foundation frontier.

\section{Fixed Inputs}

\begin{definition}[Foundation-frontier fixed inputs]
The criterion recorded here uses the following fixed inputs.
\begin{enumerate}
    \item The benchmark exact-closure package, which fixes the public one-metric GR target of the repository.
    \item The post-bridge routing record, which already moved the live burden upstream from same-package observer-side closure to substrate derivation of that benchmark package.
    \item The recorded ground-from-foundation continuation, which moved the active open burden to foundation scope without by itself settling the benchmark-facing routing question there.
    \item The ground-generating substrate foundation dynamics necessity / uniqueness theorem, which proves that the benchmark-relevant foundation core \(\FoundationCore\) is forced up to core-faithful equivalence once the fixed ground package is required.
    \item The foundation-from-root theorem, which records one deeper continuation below that foundation frontier while preserving the same benchmark package.
\end{enumerate}
\end{definition}

\begin{remark}
The present note does not reopen the mathematics of those manuscripts. It only records the routing consequence of reading them together under the active file-map safeguard against recursive depth drift.
\end{remark}

\section{Why the Foundation-Core Theorem Remains the Frontier}

\begin{definition}[Benchmark-neutral deeper relay]
A \emph{benchmark-neutral deeper relay} is a manuscript that preserves the same benchmark one-metric output, the same ordinary matter route, and the same settled observer-side closure verdict while merely relocating the unexplained substrate layer to a deeper class without providing a new compression, necessity, uniqueness, or comparably sharp routing gain.
\end{definition}

\begin{proposition}[The foundation-core compression theorem is the live frontier]
At the current foundation stage of the primitive-reservoir observer-reduction line, the ground-generating substrate foundation dynamics necessity / uniqueness theorem remains the live benchmark-facing frontier.
\end{proposition}

\begin{proof}
That theorem removes benchmark-relevant freedom at the foundation level itself: once the fixed ground package is required, the benchmark-relevant foundation core \(\FoundationCore\) is forced up to core-faithful equivalence. This is exactly the kind of necessity / uniqueness gain that changes benchmark-facing status under the routing rule. By contrast, a manuscript that only preserves the same benchmark package while pushing the unexplained stopping point deeper does not alter benchmark-facing status unless it adds a comparably sharp result. Therefore the live frontier remains the foundation-core compression theorem rather than every later deeper continuation recorded below it.
\end{proof}

\begin{corollary}[Status of the foundation-from-root theorem]
The foundation-from-root theorem is a legitimate deeper continuation on the active line, but under the current routing safeguard it is not by itself the default current benchmark-facing frontier.
\end{corollary}

\begin{proof}
The theorem records a deeper root class whose projected exact-shell output recovers the already fixed foundation package. That matters for continuity of the deeper chain. But because the induced benchmark-facing output remains the same and the benchmark-relevant internal foundation freedom had already been compressed by the preceding necessity / uniqueness theorem, the deeper relay does not on its own supersede the routing status of that frontier theorem.
\end{proof}

\section{The Next Benchmark Criterion}

\begin{definition}[Admissible next benchmark-facing foundation manuscript]
At the current frontier, a manuscript counts as an admissible next benchmark-facing foundation move only if it does at least one of the following:
\begin{enumerate}
    \item derives or justifies the ground-generating substrate foundation dynamics from still more primitive substrate structure in a way that does more than merely relocate the unexplained layer deeper, or
    \item proves a genuinely smaller collapse strictly inside the benchmark-relevant foundation core \(\FoundationCore\).
\end{enumerate}
\end{definition}

\begin{theorem}[Next honest benchmark-facing task at foundation scope]
The next honest benchmark-facing manuscript at the current foundation frontier must either derive or justify the ground-generating substrate foundation dynamics from still more primitive substrate structure in a benchmark-relevant way, or else prove a genuinely smaller collapse inside the benchmark-relevant foundation core. A manuscript that only repeats shell, lift, support, reservoir, or ground data at a deeper level without such a new gain does not satisfy the criterion.
\end{theorem}

\begin{proof}
By Proposition 1, the live frontier is already the theorem that compresses the benchmark-relevant foundation freedom. Therefore a second manuscript at foundation scope must improve on that status rather than merely accompany it. A deeper-origin manuscript can do so only if it changes benchmark-facing status by more than depth alone. If it simply reproduces the same foundation package through another relay, the routing rule classifies it as benchmark neutral. The remaining honest alternatives are exactly those stated in the definition: either a more primitive derivation that changes benchmark-facing status, or a sharper internal collapse of the current foundation core.
\end{proof}

\section{What the Next Move Should Not Be}

\begin{proposition}[Screened next moves]
At the current foundation frontier, none of the following is the default next benchmark-facing move:
\begin{enumerate}
    \item another shell, lift, support, reservoir, or ground restatement;
    \item a reopening of stationary reduction, stationary Hessians, spectral generators, or selector mechanics;
    \item a manuscript that introduces a second operative metric, enlarges the ordinary matter law, or uses stronger promotion language.
\end{enumerate}
\end{proposition}

\begin{proof}
Items (1) and (2) fail the preceding criterion because they do not directly address the current foundation-level burden. They revisit lower or screened layers rather than settling the live open issue at benchmark scope. Item (3) violates the fixed claim boundary of the benchmark package and therefore cannot count as a disciplined next step on the present line. The benchmark package remains the exact one-metric GR package already fixed by adoption, so any admissible next manuscript must preserve that boundary.
\end{proof}

\begin{remark}
This screening statement is not a denial that those lower-layer manuscripts exist or matter historically. It is a routing statement: they are not the default way to spend the next benchmark-facing move from the current foundation frontier.
\end{remark}

\section{Conclusion}

The positive result of this note is routing discipline at the current foundation frontier. The ground-generating substrate foundation dynamics necessity / uniqueness theorem remains the live benchmark-facing gain because it already spends the honest internal compression move available at the current foundation level. The deeper foundation-from-root theorem remains recorded and active, but under the current safeguard it is a deeper continuation rather than the default next benchmark-facing frontier.

The scope remains narrow and honest. The benchmark package is unchanged: the one operative metric is still exactly \(\CompletedMetric\), the ordinary matter route is unchanged, there is no second operative metric, and the low-energy matter law is not enlarged. Nothing here upgrades adoption to derivation or licenses stronger promotion language.

The next open burden is therefore clear. The next benchmark-facing manuscript should derive or justify the ground-generating substrate foundation dynamics from still more primitive substrate structure in a way that does more than another same-output relay, unless the sharper honest gain available instead is a genuinely smaller collapse inside the benchmark-relevant foundation core. That is the clean next manuscript target at current scope.

\end{document}
